\documentclass[journal,onecolumn]{IEEEtran}

\usepackage[utf8]{inputenc}
\usepackage[T1]{fontenc}
\usepackage{amsmath}
\usepackage{amssymb}
\usepackage{graphicx}
\usepackage{booktabs}
\usepackage{array}
\usepackage{stfloats}
\usepackage{url}
\usepackage{xr-hyper}
\usepackage[hidelinks]{hyperref}
\externaldocument[M-]{gvs5h_tcyb}

\newcommand{\plotdir}{figures}
\newcommand{\panelplot}[1]{%
  \includegraphics[width=\linewidth,height=0.42\textheight,keepaspectratio]{#1}}
\let\halfplot\panelplot
\renewcommand{\thefigure}{S\arabic{figure}}
\renewcommand{\thetable}{S\arabic{table}}
% The figure files ask for page tops ([!t]); here they should follow the text in order.
\usepackage{float}
\renewenvironment{figure*}[1][]{\begin{figure}[H]}{\end{figure}}
\newenvironment{chartfigure}{\begin{figure}[H]}{\end{figure}}
\renewenvironment{table*}[1][]{\begin{table}[H]}{\end{table}}

\begin{document}

\title{Supplementary Material for ``GVS5H: Zero-Shot Self-Orchestration With Ledger-Based
Control Improves Coding in Language Models''}
\author{Anonymous Authors}
\markboth{IEEE Transactions on Cybernetics}%
{Anonymous Authors: Zero-Shot Self-Orchestration With Ledger-Based Control}
\maketitle

\section{Two Corrections to the LiveCodeBench Evaluator}
\label{app:evaluator}

%% Generated by paper_plot_script/make_regrade_table.py -- do not edit.
% T-CYB copy: hand-adapted to IEEEtran (full-width floats, IEEE captions); regenerate
% from ../ and re-apply if the source changes.
\begin{table*}[!t]
\centering
\caption{Every Reported Arm as the Stock Evaluator Graded It, and After Each Correction}
\label{tab:regrade}
\setlength{\tabcolsep}{4pt}
\begin{tabular}{@{}lcccc@{\hspace{1.2em}}cccc@{}}
\toprule
& \multicolumn{4}{c}{\textbf{Single call}} & \multicolumn{4}{c}{\textbf{With manager}} \\
\cmidrule(r){2-5}\cmidrule(l){6-9}
\textbf{Model} & \textbf{\shortstack{As\\graded}} & \textbf{\shortstack{$\Delta$\\harness}}
& \textbf{\shortstack{$\Delta$\\keys}} & \textbf{\shortstack{Re-\\ported}}
& \textbf{\shortstack{As\\graded}} & \textbf{\shortstack{$\Delta$\\harness}}
& \textbf{\shortstack{$\Delta$\\keys}} & \textbf{\shortstack{Re-\\ported}} \\
\midrule
\multicolumn{9}{@{}l}{\emph{Pinned backends, five passes}} \\
Claude Fable 5 & 87.4 & $+0.0$ & $+3.0$ & \textbf{90.4} & --- & --- & --- & --- \\
GPT-5.6-Terra & 69.8 & $+8.4$ & $+2.6$ & \textbf{80.8} & 72.6 & $+12.6$ & $+2.8$ & \textbf{88.0} \\
GPT-5.6-Luna & 60.4 & $+8.0$ & $+2.0$ & \textbf{70.4} & 70.4 & $+8.4$ & $+2.4$ & \textbf{81.2} \\
Qwen3.8-27B & 62.4 & $+2.6$ & $+1.8$ & \textbf{66.8} & 84.4 & $+5.4$ & $+2.6$ & \textbf{92.4} \\
Qwen3.8-Flash-Next & 78.8 & $+2.4$ & $+3.0$ & \textbf{84.2} & 87.4 & $+2.6$ & $+3.0$ & \textbf{93.0} \\
DeepSeek-V4.1-Flash & 84.2 & $+1.2$ & $+3.0$ & \textbf{88.4} & 87.6 & $+1.4$ & $+3.0$ & \textbf{92.0} \\
\midrule
\multicolumn{9}{@{}l}{\emph{OpenRouter-served, one pass}} \\
Opus-5 & 85.0 & $+0.0$ & $+3.0$ & \textbf{88.0} & 91.0 & $+1.0$ & $+3.0$ & \textbf{95.0} \\
Kimi-K3 & 82.0 & $+1.0$ & $+3.0$ & \textbf{86.0} & 82.0 & $+1.0$ & $+3.0$ & \textbf{86.0} \\
MiniMax-M3 & 60.0 & $+0.0$ & $+2.0$ & \textbf{62.0} & 66.0 & $+0.0$ & $+3.0$ & \textbf{69.0} \\
Qwen3.6-35B & 25.0 & $+0.0$ & $+1.0$ & \textbf{26.0} & 43.0 & $+0.0$ & $+1.0$ & \textbf{44.0} \\
\bottomrule
\end{tabular}
\par\vspace{0.5ex}
\parbox{\linewidth}{\footnotesize\raggedright (Section~\ref{app:evaluator}). \textbf{As graded} is LiveCodeBench's own
verdict on the stored generations. \textbf{Harness} is the move contributed by running
each submission as a subprocess, which is what repairs the \texttt{sys.stdin} mock and
\texttt{sys.stdout.buffer}; \textbf{keys} is the move contributed by judging the four
problems that accept more than one answer by the contest's rule. \textbf{Reported} is the
number used everywhere else in this paper. Upper block: five-pass means at a 128k cap,
reasoning on (Table~\ref{M-tab:main} of the main paper). Lower block: the 128k reasoning-on, one-pass
OpenRouter arms (Table~\ref{M-tab:openrouter} of the main paper); Qwen3.5-9B is unusable in this condition and
is omitted. Across the 48 arm-passes of the two remaining OpenRouter conditions the harness
correction moves six by one point and the key correction moves each arm by 0 to 3 points.
Qwen3.8-27B's single call is the one row where \textbf{as graded} is not the stock evaluator
on the code being reported: that arm is cap-matched by replaying its 250k generations at 128k
(Appendix~\ref{M-app:protocol} of the main paper) and the truncated solutions were never put through the stock
evaluator, so the column is the 250k grading and the two $\Delta$ columns carry the cap match
as well as the corrections. Against that baseline the row gains 35 problem-passes and loses
13, eight of them generations whose 128k prefix contains no code at all.}
\end{table*}


While inspecting transcripts for Section~\ref{M-sec:discussion} of the main paper we found two defects in the grading
path: one in the LiveCodeBench harness, one in its answer keys. Every number in
Section~\ref{M-sec:results} of the main paper is reported after fixing both; no model was re-run.
Table~\ref{tab:regrade} gives every reported arm as the stock evaluator graded it and after each
correction. Neither correction is a constant offset, and the larger of the two is the answer keys
for some models and the harness for others.

\subsection{The \texttt{sys.stdin} Mock} The evaluator does not run a candidate as a subprocess.
Instead, it executes it in-process with
\texttt{sys.stdin} replaced by a mock. That mock's binary view implemented \texttt{readline()} as
\texttt{inputs.split(b"\textbackslash n")[0]}, which is a stateless expression that returns the
first line on every call. A program reading multi-line input through
\texttt{sys.stdin.buffer.readline()} therefore read line 1 repeatedly and failed however correct it
was. The mock's text-mode \texttt{sys.stdin.readline} was implemented as a proper iterator and
\texttt{buffer.read()} returns the whole payload, so neither was affected, which is why the
common \texttt{sys.stdin.buffer.read().split()} idiom never exposed it. We replaced the mock's
binary view with one backed by a \texttt{BytesIO} so that reads advance a position, and re-scored
every stored generation.

We report it explicitly for two reasons. First, it interacts badly with the v2 verifier: that step
runs the candidate as a real subprocess (step 5 of Section~\ref{M-sec:method} of the main paper), where
\texttt{buffer.readline()} works correctly, so the manager was told its solution passed the public
samples for programs the grader then marked wrong - the loop's one external signal certifying a
wrong answer. Second, the exposure is uneven across models rather than a constant offset: 371 of the
5{,}425 pinned-backend outputs containing code use the idiom, and 86\% of those were scored wrong
against 18\% for the outputs that do not. Fable 5 never uses it and its scores are unchanged to the
decimal. The OpenRouter-served set reaches for it once in 4{,}764 generations containing code, and
then through an alias - Kimi-K3's single call on \texttt{abc396\_g} binds
\texttt{sys.stdin.buffer} to a name and calls \texttt{readline()} on that. Every OpenRouter-served
number is re-scored all the same, and the repaired mock moves one cell: Kimi-K3's 128k
reasoning-on single arm, from 82 to 83 (Table~\ref{tab:regrade}, where that cell reaches its
reported 86 through the answer-key correction below). Because the idiom is a matter of coding
style, a harness bug of this shape is not a wash across a leaderboard - it penalizes the models
whose style trips it, and its $\Delta$ harness column ranges from $+0.0$ for Fable 5, which
never writes it, to $+12.6$ for GPT-5.6-Terra's manager arm.

\subsection{Answer Keys That Accept One of Many Valid Answers} LiveCodeBench grades by exact
string match against a single stored output. Four LCB-100 problems accept more than one:
\texttt{abc396\_e}, \texttt{arc190\_a} and \texttt{arc195\_c} ask for any construction that
attains the optimum, and \texttt{abc385\_f} asks for a real number, which its contest accepts
within a relative $10^{-6}$. We judge those four by the rule their contest used. Each of the three
constructive problems gets a checker that parses the submitted answer, tests it against the
problem's own constraints, and requires it to attain the objective value of the stored reference:
for \texttt{abc396\_e}, that the sequence is non-negative, satisfies every XOR constraint and sums
to what the reference sums to; for \texttt{arc190\_a}, that the chosen operations cover $1..N$ and
number what the reference says is optimal; for \texttt{arc195\_c}, that the Yes/No verdict matches
the reference and that any placement offered uses the requested counts of each colour, on distinct
cells, with every consecutive pair a step of the kind its colour allows. \texttt{abc385\_f} is
compared token by token at $10^{-6}$ relative tolerance. This class of defect is one LiveCodeBench documents upstream - its errata
list eleven earlier problems whose keys accept multiple answers, none of them in LCB-100 - and
it is a property of the benchmark, not of any model, so it applies to every arm of both sets and
moves 104 of their 118 arm-passes. Re-judging
runs each submission as a real subprocess rather than in-process, which is also what repairs the
mock above and makes \texttt{sys.stdout.buffer} write as it does on the contest judge.

Both corrections are conservative in the same two ways. They re-run only the submissions the stock
evaluator marked wrong, so neither can turn a pass into a failure and the corrections move every
arm of Table~\ref{tab:regrade} upward; and among those they re-run only
the ones with a known reason to be misjudged - a failure on one of the four problems above, or
code that touches \texttt{sys.stdin.buffer} or \texttt{sys.stdout.buffer}. A submission that the
in-process evaluator failed for some other reason keeps that verdict. Re-executing every stored
generation instead, passing and failing alike, finds 25 further submissions across 11{,}800
records that pass every hidden test as a subprocess; they are spread over 22 arms, at most three
to an arm, and we leave them as the stock evaluator scored them.

\section{Experimental Conditions}
\label{sec:conditions}

The pinned-backend set runs at one setting: 128k output cap, reasoning on,
$\times 5$ passes, with the five models in both arms and Fable 5 single-only. The
OpenRouter-served set runs at three: 128k reasoning-on $\times 1$; 16k reasoning-off
$\times 5$, whose repeated passes give run-to-run variance and paired tests; and 128k reasoning-off
$\times 1$, the same without the small cap, isolating it as a confound.

\section{Cost per Arm}

% Generated by paper_plot_script/make_figures_tex.py -- do not edit.
% T-CYB copy: hand-adapted to IEEEtran (full-width floats, IEEE captions); regenerate
% from ../ and re-apply if the source changes.
% The caption text and the table numbers live in that chart's plot script;
% re-run the script to update this file.
\begin{chartfigure}
\centering
\halfplot{\plotdir/what_one_pass_costs_lcb-100_5_passes_single_call_vs_manager_against_fable_5_light.pdf}
\caption{The scaffold's bill. Cost of one pass over the same 100 problems; Table~\ref{M-tab:cost} of the main paper is the arithmetic behind every bar. No cached-input discount is taken, and Qwen3.8-27B is priced at OpenRouter market rates. The top bar of each pair is the single call, light; the manager is under it, dark; all are at a 128k output cap, and colour is the model. Fable 5 is first; the models below it run top to bottom from the cheapest single call to the dearest. Table~\ref{M-tab:cost} carries the tests. Unbracketed comparisons: every manager $-$ single increase is significant, and GPT-5.6-Luna's manager undercuts GPT-5.6-Terra's single call by \$1.91. At 1.28x the assumed Qwen rate (\$0.273/\$3.26 per MTok, inside the spread across hosted providers) the manager's cost advantage over Fable 5 disappears entirely - a larger uncertainty than any p-value here.}
\label{fig:cost}
\end{chartfigure}


\section{Accuracy by Model Scale}

Figs.~\ref{fig:scale-on} and~\ref{fig:scale-off} plot the one-pass OpenRouter-served results of
Fig.~\ref{M-fig:scale-bars} of the main paper against model parameter count; the five-pass
condition is Fig.~\ref{M-fig:scale-16k} of the main paper.

% Generated by paper_plot_script/make_figures_tex.py -- do not edit.
% T-CYB copy: hand-adapted to IEEEtran (full-width floats, IEEE captions); regenerate
% from ../ and re-apply if the source changes.
% The caption text lives in that chart's plot script, next to the numbers it
% describes; re-run the script to update this file.
\begin{figure*}[!t]
\centering
\panelplot{\plotdir/manager_vs_single_call_across_model_scale_lcb-100_1_pass_128k_max_tokens_reasoning_on_light.pdf}
\caption{Accuracy by model scale, reasoning on, 128k, one pass. The same numbers as Fig.~\ref{M-fig:scale-bars}(a), against parameter count. Dashed = single call, solid = with manager; the block under each tick gives that model's $\Delta$ and Tukey group, and the legend each arm's gain from the smallest model to the largest. }
\label{fig:scale-on}
\end{figure*}


% Generated by paper_plot_script/make_figures_tex.py -- do not edit.
% T-CYB copy: hand-adapted to IEEEtran (full-width floats, IEEE captions); regenerate
% from ../ and re-apply if the source changes.
% The caption text lives in that chart's plot script, next to the numbers it
% describes; re-run the script to update this file.
\begin{figure*}[!t]
\centering
\panelplot{\plotdir/manager_vs_single_call_across_model_scale_lcb-100_1_pass_128k_max_tokens_reasoning_off_light.pdf}
\caption{Accuracy by model scale, reasoning off, 128k, one-pass control. The same numbers as Fig.~\ref{M-fig:scale-bars}(c), against parameter count. Dashed = single call, solid = with manager; the block under each tick gives that model's $\Delta$ and Tukey group, and the legend each arm's gain from the smallest model to the largest. }
\label{fig:scale-off}
\end{figure*}


\end{document}
